\documentclass[11pt]{article}
\usepackage[T1]{fontenc}
\usepackage{lmodern}
\usepackage[margin=1in]{geometry}
\usepackage{amsmath,amssymb,amsthm,mathtools}
\usepackage{booktabs,longtable,array}
\usepackage{microtype}
\usepackage{xcolor}
\usepackage{enumitem}
\usepackage{xurl}
\usepackage[colorlinks=true,linkcolor=blue!45!black,urlcolor=blue!45!black,citecolor=blue!45!black]{hyperref}
\hypersetup{pdftitle={Twelve positive unknowns for the consistency of ZFC},
 pdfauthor={Prepared for Vladimir Reshetnikov},
 pdfsubject={Explicit Diophantine constructions in 12 positive or 26 integer unknowns}}
\usepackage{fancyhdr}
\pagestyle{fancy}\fancyhf{}
\fancyhead[L]{\small Diophantine forms of $\operatorname{Con}(\mathrm{ZFC})$}
\fancyhead[R]{\small Twelve positive unknowns}
\fancyfoot[C]{\thepage}
\setlength{\headheight}{14pt}
\setlength{\emergencystretch}{3em}
\setlist{itemsep=3pt,topsep=5pt}
\newcommand{\Z}{\mathbb Z}
\newcommand{\N}{\mathbb N}
\newcommand{\Pos}{\mathbb Z_{>0}}
\newcommand{\Con}{\operatorname{Con}}
\newcommand{\pc}{\operatorname{pc}}
\newcommand{\ZFCRef}{\operatorname{Ref}_{\mathrm{ZFC}}}
\newcommand{\Norm}{\mathcal R}
\newcommand{\src}[1]{\nolinkurl{#1}}
\newtheorem{theorem}{Theorem}[section]
\newtheorem{lemma}[theorem]{Lemma}
\newtheorem{proposition}[theorem]{Proposition}
\newtheorem{corollary}[theorem]{Corollary}
\theoremstyle{remark}\newtheorem{remark}[theorem]{Remark}
\title{\textbf{Twelve positive unknowns for the consistency of ZFC}\\[6pt]
\large A separating root construction, 26 integer unknowns,\\and a smaller finite coefficient compiler}
\author{Prepared for Vladimir Reshetnikov\\\small AI-assisted mathematical research synthesis}
\date{22 October 2026}
\begin{document}
\maketitle
\begin{abstract}
We combine the current ProveIt blank-machine Diophantine constructions and
give further reductions in three measures. First, an explicitly specified polynomial in
\textbf{12 positive integer unknowns}, of exact total degree \textbf{1369},
has a fully charged \textbf{144-operation} evaluator. Its solvability is
equivalent, by ordinary mathematics formalizable in ZFC, to a finite ZFC
refutation. Its three-equation version has maximum degree 1360 and uses
140 operations. The new step combines four auxiliary witnesses into one
weighted integer. A divisibility adjustment to a Pell multiplier makes the
reconstructed rational witnesses integers; this is essential to soundness.
Second, robust separating weights give \textbf{26 unrestricted integer
unknowns}, degree 3560 and 199 operations; an independent signed-scale
argument gives lower-degree alternatives in \textbf{27 integer unknowns}.
Third, a refined local alphabet
and compressed layout reduce the distinguished coefficient from
$B=2^{5^{18}}$ to \textbf{$B=2^{5^{16}}$}; the new alphabet has 5263 tiles.
The archive contains the complete arithmetic circuits, coefficient recipe,
finite verification programs, source provenance, and proof notes.

The comparison is pinned to ProveIt commit
\src{f0ae4e0991ca44a7bb4e5671d42e33d775e48035}.
These are reductions in specified measures, not a proof of an absolute
minimum, a new universal-pair theorem, or a proof of ZFC consistency.
The new Diophantine proofs are not kernel checked. The inherited exact-machine
endpoint has Lean and Rocq evidence, which is kept distinct from the new work.
\end{abstract}
\clearpage
\tableofcontents
\clearpage

\section{Outcome, comparison, and conventions}\label{sec:results}

There is no single ordering of the constructions by simplicity. The fewest
unknowns, lowest degree, shortest arithmetic evaluator, and smallest fixed
coefficients favor different presentations. The main new result is a further
reduction in the number of positive unknowns, together with independent
improvements for unrestricted integers and coefficient size.

\begin{theorem}[Fixed-coefficient consistency equation]\label{thm:main}
Bind the four coefficient ports $(\beta,\kappa,M_C,M_F)$ to the complete
5263-tile recipe of Section~\ref{sec:coefficients}. Let $P_{12}$ be the
ordinary integer polynomial defined in Section~\ref{sec:twelve}. Then
\[
 \mathrm{ZFC}\vdash
 \left(\exists\mathbf x\in\Pos^{12}\;P_{12}(\mathbf x)=0\right)
 \longleftrightarrow \ZFCRef.
\]
Consequently,
\[
 \boxed{\quad\mathrm{ZFC}\vdash
 \left(\forall\mathbf x\in\Pos^{12}\;P_{12}(\mathbf x)\ne0\right)
 \longleftrightarrow\Con(\mathrm{ZFC}).\quad}
\]
The same existence equivalence holds for the stated three-equation system
and for each of the five 27-integer constructions in Section~\ref{sec:integer},
as well as the 26-integer constructions in Section~\ref{sec:integer26}.
\end{theorem}

Here $\ZFCRef$ means existence of a finite syntactic proof of contradiction from
an ordinary effective presentation of ZFC, and $\Con(\mathrm{ZFC})$ is its
negation. Section~\ref{sec:logic} explains the imported calibration and the
formalization boundary.

\subsection{The current repository frontier}

The supplied collection now contains seven wave-2 packages, numbered 09--15.
They supersede eight earlier packages. Its latest audited reconciliation
explicitly distinguishes several coincident discoveries and corrects earlier
degree, domain, and coefficient-binding descriptions~\cite{catalogue}.
We compare against that wave-2 tree, not the superseded wave-1 frontier.

\begin{center}\small
\begin{tabular}{>{\raggedright\arraybackslash}p{5.2cm}rrrr}
\toprule
Construction & Unknowns & Equations & Degree & Ops\\
\midrule
Existing minimum-operation product & $15\Pos$ &1&155&67\\
Existing ordinary system (package 12)&$15\Pos$&6&28&60\\
Existing lowest-degree 15-variable system&$15\Pos$&6&22&61\\
Existing weighted root product (09)&$13\Pos$&1&515&85\\
Existing shorter root product (09)&$13\Pos$&1&571&84\\
Existing quartic lift (09,12)&$32\Pos$&1&4&---\\
Existing integer ordinary product (14)&$29\Z$&1&188&121\\
\midrule
\textbf{New separating-root product}&$\mathbf{12\Pos}$&1&1369&144\\
\textbf{New separating-root system}&$\mathbf{12\Pos}$&3&1360&140\\
New 12-variable sum of squares&$12\Pos$&1&2720&147\\
\textbf{New robust integer product}&$\mathbf{26\Z}$&1&3560&199\\
New robust integer system&$26\Z$&3&3544&195\\
New robust integer sum of squares&$26\Z$&1&7088&202\\
\textbf{New integer ordinary product}&$\mathbf{27\Z}$&1&172&116\\
New integer ordinary system&$27\Z$&6&46&107\\
New integer ordinary sum of squares&$27\Z$&1&92&124\\
\bottomrule
\end{tabular}
\end{center}

The symbol $n\Pos$ means $n$ strictly positive integer unknowns; $n\Z$ means
$n$ unrestricted integer unknowns. For a system, degree is the maximum degree
of its residuals after all abbreviations are substituted. The comparison
counts every binary addition, subtraction, and multiplication, including
multiplication by a fixed numeral. Register reuse, constants, and final
comparisons are free. If a system is required to output residuals rather than
compare its expressions to constants, the new 12-variable residual evaluator
costs 142 operations instead of 140. This distinction is recorded explicitly.

The 12-variable result improves the reviewed explicit positive-variable
frontier, but its degree and evaluator are larger than those of the
13-variable products. The 26-integer version further reduces the integer
coordinate count, at a substantial degree and operation cost.
The 27-integer ordinary product improves the previous
29-integer ordinary product in all three displayed numerical measures.
Neither change improves the 67-operation positive-variable evaluator.

\begin{center}\small
\begin{tabular}{lrrr}
\toprule
Coefficient backend & Local tiles & $B$ & Bits of $B-1$\\
\midrule
Existing T4 (11)&37075&$2^{5^{18}}$&3814697265625\\
New refined T4&5263&$2^{5^{16}}$&152587890625\\
\bottomrule
\end{tabular}
\end{center}

The new base coefficient has exactly 25 times fewer bits. It is still far too
large for a convenient expanded numeral: the base alone takes about 19.1 GB
in binary. We therefore specify the four integers by exact terminating finite
recipes and bind each arithmetic circuit to the whole tuple. A compact recipe
does not make these coefficients numerically small.

Classical nine-variable constructions over the nonnegative integers exist
(and translate to positive integers), at much larger general degree
bounds, and relation combining is a classical technique~\cite{mr75,jones82,bayer}.
The new result concerns this explicitly bound K215 compiler and the displayed
circuit. It is not a claim that twelve is the minimum possible number of
unknowns, nor that $(12,1369)$ is a universal pair for all Diophantine sets.

\section{The inherited arithmetic interface}\label{sec:interface}

The coefficient contract uses $\beta=B-1$, a transport coefficient $\kappa$,
a low mask $M_C$, and a \emph{shifted} high mask $M_F$. If the native high mask
is $M_F^{\rm nat}$, then $M_F=M_F^{\rm nat}+B-1$. In particular
\begin{gather}
 B\ge16,\quad 0<M_C<B-1,\quad0<M_F^{\rm nat}<B-1,\notag\\
 M_C\equiv2\pmod4,\quad M_F^{\rm nat}\equiv4\pmod8,\quad
 \pc(M_C)+\pc(M_F^{\rm nat})=\log_2B.\label{eq:maskcontract}
\end{gather}
Those arithmetic properties accompany the full local-tableau, carry,
synchronization, and canonical-witness contract. They alone would not be
sufficient to encode a machine.

With the common positive coordinates, put
\begin{align}
q&=\beta J+1,& C&=q-F-\alpha,& X&=wq,&Y&=sq^3,\notag\\
E&=XY,& k&=\eta+\zeta,&c&=kY+\eta,&a&=E+Y,\notag\\
H&=4a+3,&\Delta&=a^2+H=(a+2)^2-1,\label{eq:prefix}\\
P&=X+ac,&S&=1+\Delta c^2,& L&=EkY,&A&=1+L(L+k).
\notag
\end{align}
Define the actual packed rank and the two retained factors by
\begin{align}
 R&=\bigl((q-1)(q-F)+\alpha\bigr)(q^2-1)+(M_C+qM_F)J,
 \label{eq:rank}\\
 N_k&=k-hE-R,&
 N_t&=(\kappa+w)C+(q-F)-z(q-1).\label{eq:units}
\end{align}
All capitalized expressions are abbreviations for substituted polynomials,
not additional unknowns.

The earlier normalized construction uses fifteen positive coordinates
\[
 J,F,\alpha,z,f,h,i_0,T,s,w,\tau,\eta,\zeta,y,\gamma.
\]
Its system can be written
\begin{gather}
 \tau^2=A,\qquad (P+H\gamma)^2=S,\qquad
 f^2=1+\Delta i_0^2c^4,\notag\\
 Q_0\bigl(c(Tf-1)-Rf^2\bigr)^2-(Q_0-1)y^2=1,
 \qquad Q_0=(\Delta i_0c^2)^2,\label{eq:parent}\\
 N_k=N_t=1.\notag
\end{gather}
The inherited theorem identifies existence of such a tuple with a blank
halting run of the literal K215 machine~\cite{blank,rank,signed29}.

\subsection{What is used from the signed theorem}

The soundness argument also permits $f,T,\tau,y,\gamma,h$ to be signed,
provided the compiler coordinates $J,F,\alpha,z,s,w,\eta,\zeta$ are positive
and the normalized multiplier is nonzero. The direct nine-positive, six-signed theorem is proved in package 11's
\src{reduction/SIGNED_DOMAIN_PROOF.md}, copied here as
\src{upstream/signed_domain_core.md}. It also proves the ordinary-strong
version. Its hypothesis is met by the positive multiplier used below.
The related 29-integer theorem in package 14 adds oddness and a square
scale for a different parametrization; those extra restrictions are not
needed for the 12-variable reduction.

For the unit-product version, it is enough initially to have
$N_k=N_t=\varepsilon\in\{1,-1\}$, with the four norm equations equal to one.
The proof establishes $\varepsilon=1$ after recovering the Pell indices.
In particular it does not assume the index or transport sign before the
required witnesses have been restored.

Here are the main logical steps of that imported arithmetic proof. Write
$\chi_U(n)+\psi_U(n)\sqrt{U^2-1}=(U+\sqrt{U^2-1})^n$ for integer $U\ge2$.
The first and main equations classify the relevant roots as
$k=2\psi_{2XY^2+1}(n)$ and $c=\psi_{a+2}(p)$. The index equation bounds $n$;
the main norm and normalized strong divisibility give a larger index $m$
with $pc\mid m$. The two literal auxiliary congruences, modulo $f$ and $c$,
give $R\equiv\pm p\pmod c$. Growth bounds force $p=R$.
The first-index wrap and the negative sign are then excluded.
The lower Pell-ratio estimate and the congruence modulo $H$ give $X=2^R$.
Half-binomial extraction and population inversion recover the two native
masks, and the coefficient compiler decodes an actual finite computation.
The signed main root is controlled by the separate signed-$\gamma$ argument:
its wrong sign would give $X2^R\equiv-1\pmod H$ in a strict interval too small
to permit that congruence. These statements, their hypotheses, and the
supporting source proofs are included in \src{upstream/}.

\subsection{Bounds available before restoring any root}

Suppose only that $N_t$ is a unit. Since
\[
 N_t=(\kappa+w)C+1-F-(z-1)(q-1),
\]
one has $C\ge0$. The boundary $C=0,F=2,z=1$ can yield $N_t=-1$,
so strict positivity cannot yet be used. The authentic mask bounds give
\begin{gather}
 q(q^2-1)<R<q^4-q^3+q^2-1<q^4-2,\notag\\
 q\ge16,\qquad X\ge q,\qquad Y\ge q^3\ge4096,
 \qquad E\ge q^4>R+2.\label{eq:weak}
\end{gather}
Also $kY<c<k(Y+1)<2kY$, $c>1$, and $a+2<2E$.
These bounds use no eliminated root. That fact makes the proof in the next
section noncircular.

\section{A polynomial in twelve positive unknowns}\label{sec:twelve}

The supplied coordinates are exactly
\[
 \boxed{J,F,\alpha,z,h,i,T,s,w,\eta,\zeta,t\in\Pos.}
\]
The four former coordinates $\gamma,\tau,f,y$ are replaced by one coordinate
$t$. The polynomials in \eqref{eq:prefix}--\eqref{eq:units} are retained.

\subsection{The divisibility adjustment and the separating weights}

Define
\begin{align}
 G&=Hic^2,&D&=1+\Delta G^2,& Q&=\Delta(D-1)=(\Delta G)^2,
 &d&=Q-1,\notag\\
 v&=cT,&b&=c+RD,&e&=Q(Dv^2+b^2)-1,&B_0&=2Qvb.
 \label{eq:aux12}
\end{align}
The four equations whose witnesses will be eliminated are
\begin{equation}
 (P+H\gamma)^2=S,\qquad \tau^2=A,\qquad f^2=D,
 \qquad dy^2=e-B_0f.\label{eq:fourroots}
\end{equation}
The last equality is equivalent, when $f^2=D$, to
\[
 Q(vf-b)^2-(Q-1)y^2=1.
\]
This is the original normalized system with $i_0=Hi$, since
$vf-b=c(Tf-1)-Rf^2$. The factor $H$ in $G$ is the new divisibility
adjustment. Its role is the exact identity
\begin{equation}
 Q\equiv0\pmod{H^2},\qquad \gcd(H,d)=1.\label{eq:coprime}
\end{equation}
Completeness of this restriction will be proved below.

Use the positive polynomial weights
\begin{align}
 M&=1+8THic^2+16REH^2i^2c^3,\notag\\
 N&=1+2HM.\label{eq:weights12}
\end{align}
The coordinate $t$ is intended to encode
\begin{equation}
 t=y+M\tau+N\gamma.\label{eq:trace}
\end{equation}
In particular $N\equiv1\pmod H$. Neither weight is a new unknown.

\subsection{An explicit formula using only polynomial operations}

Set $W=Ht+NP$ and $m=HM$. Define the scalar polynomials
\begin{align}
 a_0&=d(W^2+m^2A+N^2S)-H^2e,\notag\\
 a_1&=-2dWm,&a_2&=-2dWN,&a_3&=2dmN,\label{eq:normcoeff1}\\
 f_0&=a_0^2+Aa_1^2+S(a_2^2+Aa_3^2)-H^4B_0^2D,\notag\\
 f_1&=2(a_0a_1+Sa_2a_3),\notag\\
 f_2&=2(a_0a_2+Aa_1a_3),\notag\\
 f_3&=2(a_0a_3+a_1a_2).\label{eq:normcoeff2}
\end{align}
Finally put
\begin{align}
 U&=f_0^2+Af_1^2-S(f_2^2+Af_3^2),\notag\\
 V&=2(f_0f_1-Sf_2f_3),\notag\\
 \Norm&=U^2-AV^2.\label{eq:norm12}
\end{align}
The requested single equation is
\begin{equation}
 \boxed{\quad P_{12}=N_kN_t(1+\Norm)-1=0.\quad}\label{eq:final12}
\end{equation}
An equivalent existence condition is the particularly transparent system
\begin{equation}
 \boxed{\quad N_k=1,\qquad N_t=1,\qquad \Norm=0.\quad}
 \label{eq:system12}
\end{equation}
Every formula here is an ordinary polynomial with integer coefficients.
The formulas contain no variable exponent, division, congruence test,
root operation, or hidden quantified parameter. The symbols used in the
proof for eliminated roots do not occur as unknowns in the final equation.
The emitted shared circuit computes $-f_1$ and $-f_2$
to avoid unnecessary negation gates, and factors
$H^4B_0^2D$ as $4(H^2Q)^2(Dv^2)b^2$.

The construction is two successive quadratic norms. In the algebra with
formal symbols $r,u$ satisfying $r^2=A$ and $u^2=S$, the expression
\[
 \bigl[d(W-mr-Nu)^2-H^2e\bigr]^2-H^4B_0^2D
\]
reduces to $f_0+f_1r+f_2u+f_3ru$. Taking its norm first in $u$ and then in
$r$ gives \eqref{eq:norm12}. This explanation is not part of the evaluator;
the scalar equations above define the polynomial directly.

\subsection{Separation of all algebraic branches}

The first task is to show that an integer root of the eliminated equation
does not merely encode a collection of irrational or rational nonintegral
witnesses. We prove separation before any Pell root is reconstructed.

\begin{lemma}[Separating inequalities]\label{lem:separation}
Under the positive input assumptions and the weak bounds
\eqref{eq:weak}, every algebraic branch of \eqref{eq:fourroots} satisfies
\[
 |y|<M\sqrt A,\qquad S>A,\qquad
 \frac{N\sqrt S}{H}>M\sqrt A+|y|.
\]
The absolute value is the complex absolute value; the first and last
inequalities include nonreal auxiliary branches.
\end{lemma}
\begin{proof}
All of $G,D,Q,v,b$ are positive and $Q>1$. Since $G>1$,
\[
 \sqrt D<(a+2)G<2EG.
\]
Put $Z_0=v\sqrt D+b>1$. For either sign of $f=\pm\sqrt D$,
\[
 |y|^2
 \le\frac{QZ_0^2+1}{Q-1}
 \le2Z_0^2+1<4Z_0^2.
\]
Thus $|y|<2(v\sqrt D+b)$. The right side is bounded above by
\[
 4cTEG+2c+8RE^2G^2.
\]
On the other hand, $\sqrt A>L=EkY>Ec/2$. Division by this positive
quantity gives
\begin{align*}
 \frac{|y|}{\sqrt A}
 &<8TG+\frac4E+16RE\frac{G^2}{c}\\
 &=8THic^2+\frac4E+16REH^2i^2c^3<M,
\end{align*}
since $E\ge16^4>4$.

For the second claim, $c>kY$ and $\Delta-E^2>2EY$ imply
\[
 \Delta c^2>k^2Y^2\Delta
 >k^2(E^2Y^2+EY)=A-1.
\]
Here $2EY^3>EY$, as $Y\ge1$. Consequently $S>A$. Finally
\[
 \frac{N\sqrt S}{H}>2M\sqrt S>2M\sqrt A>M\sqrt A+|y|.
\]
\end{proof}

\begin{lemma}[Injective weighted encoding]\label{lem:injective}
The map $(\gamma,\tau,f,y)\mapsto y+M\tau+N\gamma$ is injective
on the finite algebraic solution set of \eqref{eq:fourroots}.
\end{lemma}
\begin{proof}
If two tuples have different $\gamma$, the absolute difference of their
$N\gamma$ terms is $2N\sqrt S/H$. By Lemma~\ref{lem:separation}, this
strictly exceeds $2M\sqrt A+|y_1|+|y_2|$, which bounds the remaining
possible cancellation. If $\gamma$ agrees but $\tau$ differs, the
$M\tau$ difference has magnitude $2M\sqrt A>|y_1|+|y_2|$.
If both agree and the encoded value agrees, then $y$ agrees. The equality
$dy^2=e-B_0f$, with $B_0>0$, forces $f$ to agree as well.
\end{proof}

\subsection{Rationality and the integrality repair}

Over an algebraic closure, with multiplicities, the exact factorization is
\begin{equation}
 \Norm=d^8H^{16}
 \prod_{(\gamma,\tau,f,y)}
 \bigl(t-y-M\tau-N\gamma\bigr),\label{eq:factorization}
\end{equation}
where the sixteen branches are those of \eqref{eq:fourroots}.
Indeed, for fixed $r=\tau$ and $u=P+H\gamma$, the inner polynomial first
factors over the two values of $f$ and then over the two values of $y$;
each of its four linear factors contributes the indicated power of $H$,
and its leading coefficient in $t$ is $d^2H^4$. Multiplying over the four
choices of $r,u$ gives \eqref{eq:factorization}.

Suppose $\Norm=0$ at integer inputs satisfying the weak bounds. As $d,H$
are nonzero, one of the algebraic tuples in \eqref{eq:factorization} has
the specified integer value $t$. Choose a finite normal extension of
$\mathbb Q$ containing all its coordinates. Every automorphism over
$\mathbb Q$ sends it to another allowed tuple and fixes its encoded value.
Lemma~\ref{lem:injective} forces the whole tuple to be fixed. The fixed
field theorem therefore makes all four coordinates rational. This argument
does not assume that the square roots generate independent extensions.

Rationality is only the first half of the argument. Since
\[
 \tau^2=A,\qquad f^2=D,\qquad (P+H\gamma)^2=S
\]
are monic square equations with integer right sides, $\tau,f,u=P+H\gamma$
are integers. Hence the reduced denominator of $\gamma$ divides $H$.
If $y=r_0/s_0$ in lowest terms, then $dy^2=e-B_0f\in\Z$ implies
$s_0^2\mid d$. Thus the denominators of $\gamma$ and $y$ are coprime by
\eqref{eq:coprime}. But
\[
 N\gamma+y=t-M\tau\in\Z,\qquad \gcd(N,H)=1.
\]
Multiplication by $N$ cannot cancel a denominator dividing $H$, and the
remaining two coprime denominators cannot cancel each other. Both therefore
equal one. All four reconstructed witnesses are integers.

This last step is the reason for strengthening the Pell multiplier by $H$.
A bare resultant, or a rationality argument without this denominator
separation, would leave a real soundness gap.

\subsection{Soundness of the system and the single equation}

For \eqref{eq:system12}, the weak bootstrap holds directly. The preceding
argument restores \eqref{eq:fourroots}, and
\[
 vf-b=c(Tf-1)-Rf^2.
\]
The ordinary strong multiplier
\[
 i_{\rm ord}=\Delta Hi>0
\]
satisfies
\[
 (i_{\rm ord}c^2)^2=\Delta(f^2-1)=Q.
\]
The reconstructed tuple therefore satisfies the signed ordinary-strong
interface exactly. Its positive ports remain positive; $f,\tau,y,\gamma$
may be signed, as permitted by the direct signed theorem in
\src{upstream/signed_domain_core.md}. The witnesses $\tau,y$ cannot be zero:
$A>1$, while $y=0$ would give $Q(vf-b)^2=1$ with $Q>1$.
The inherited rank, extraction, and compiler proof gives a blank K215 halt.

For \eqref{eq:final12}, all three integer factors are units. By
\eqref{eq:norm12}, $V$ is even and
\[
 \Norm\equiv U^2\pmod4.
\]
Thus $1+\Norm$ cannot equal $-1$, whose residue is 3. It equals 1, so
$\Norm=0$ and $N_k=N_t=\varepsilon$ for a common sign. The weak bootstrap
was explicitly proved for these unit signs; the reconstruction consequently
remains valid. Applying the signed parent proof resolves
$\varepsilon=1$ and yields the same blank halt. No sign assumption was
silently used in the root elimination.

\subsection{Completeness of the strengthened multiplier}\label{sec:complete12}

Start from a canonical accepting blank tableau and its original first/main
Pell witnesses. Its actual rank satisfies $R\equiv3\pmod4$, and
$c=\psi_{a+2}(R)$. Put $A_0=a+2$. Consider the matrix
\[
 \mathcal M=\begin{pmatrix}A_0&\Delta\\1&A_0\end{pmatrix}.
\]
Its determinant is $A_0^2-\Delta=1$. Therefore its powers in the finite
group of invertible matrices modulo $Hc^2$ are periodic. Choose $m_0>0$
to be a sufficiently large common multiple of this period and $R$.
Then $\psi_{A_0}(m_0)$ is a positive multiple of $Hc^2$. Set
\begin{equation}
 i=\frac{\psi_{A_0}(m_0)}{Hc^2},\qquad
 f=\chi_{A_0}(m_0),\qquad
 \rho=\Delta\psi_{A_0}(m_0).\label{eq:periodic}
\end{equation}
All are positive integers and $\rho^2=Q$.

To construct the auxiliary witness explicitly, use
\begin{equation}
 V_* =\frac{\chi_\rho(R)}{\rho},\qquad y=\psi_\rho(R).
 \label{eq:auxcomplete}
\end{equation}
Because $R$ is odd, $\chi_x(R)$ is divisible by $x$ as an integer polynomial.
Write $R=2j+1$ and $\chi_x(R)=xQ_j(x^2)$. The inherited elementary
Chebyshev identities are
\[
 Q_j(0)=(-1)^jR,\qquad
 Q_j(1-A_0^2)=(-1)^j\psi_{A_0}(R).
\]
Now $j$ is odd. Moreover $\rho^2\equiv0\pmod c$ and
$\rho^2\equiv1-A_0^2\pmod f$. Consequently
\begin{equation}
 V_*\equiv-R\pmod c,\qquad V_*\equiv-c\pmod f.\label{eq:quotients}
\end{equation}
The quantities
\[
 o=\frac{V_*+c}{f},\qquad T=\frac{o+Rf}{c}
\]
are positive integers. To verify the second integrality claim, note that
$of=V_*+c\equiv-R\pmod c$, $f^2\equiv1\pmod c$, and $\gcd(f,c)=1$.
Hence $(o+Rf)f\equiv0\pmod c$, so $c\mid o+Rf$.
Substitution gives
\[
 c(Tf-1)-Rf^2=V_*.
\]
The Pell identity for \eqref{eq:auxcomplete} gives
$Q V_*^2-(Q-1)y^2=1$. All other canonical positive coordinates are
unchanged. Thus the stronger system has a positive witness whenever the
machine halts. No parity restriction on $m_0$ or on the new $i$ was needed.

Finally take $t=y+M\tau+N\gamma>0$ using the canonical positive roots.
One factor of \eqref{eq:factorization} vanishes, so both the system and
the single equation have a positive solution. Together with soundness and
the machine calibration, this proves Theorem~\ref{thm:main}.

\subsection{Exact degree and arithmetic size}

All supplied unknowns have degree one. The relevant degrees are
\begin{center}\small
\begin{tabular}{lrrrrrrrrrrrr}
\toprule
Expression &$c$&$E,a,H$&$\Delta$&$R$&$A,S$&$G$&$D$&$Q,d$&$b$&$M$&$N$\\
\midrule
Degree &5&6&12&4&22&17&46&58&50&39&45\\
\bottomrule
\end{tabular}
\end{center}
The exact total degree of $\Norm$ is 1360; its degree in $t$ is 16.
Here is a uniform certificate, rather than a purely syntactic degree bound.

Scale all twelve unknowns by a formal parameter, keeping the fixed
coefficients constant. The leading homogeneous forms satisfy
\[
 A_{\rm top}=S_{\rm top}=E_{\rm top}^{,2}c_{\rm top}^{,2}.
\]
The branches of $M\tau$ have leading terms $\pm16b_{\rm top}$ of degree 50.
The quadratic equation for $\gamma$ is
\[
 H^2\gamma^2+2HP\gamma+(P^2-S)=0,
\]
where the three coefficient degrees are 12, 17, and 16; the cancellation
\[
 P^2-S=X^2+2Xac-Hc^2-1
\]
is essential. Its branches have orders 5 and $-1$ in the scaling parameter.
The large branch contributes $-64b_{\rm top}$ to $N\gamma$ at degree 50;
the small branch contributes only degree 44. Each $y$ branch contributes
$\pm b_{\rm top}$ at degree 50, independently of the leading choice of $f$.
Thus the sixteen linear factors in \eqref{eq:factorization} have nonzero
leading magnitudes
\[
 15,17,15,17,15,17,15,17,
 47,49,79,81,47,49,79,81
\]
times $b_{\rm top}$. Equivalently,
\begin{equation}
 \Norm_{1360}=C_0d_{\rm top}^{,8}H_{\rm top}^{,16}b_{\rm top}^{,16},
 \qquad C_0=15^4 17^4(47\cdot49\cdot79\cdot81)^2>0.
 \label{eq:leader}
\end{equation}
The degree is $8\cdot58+16\cdot6+16\cdot50=1360$.
This may be interpreted as a formal Laurent-series calculation over the
field of rational functions in the unscaled coordinates. No numerical
approximation to an algebraic root is being used. All leading factors are
nonzero for every valid fixed coefficient binding with $\beta>0$.

The leading forms of $N_k,N_t$ have degrees 7 and 2 and are nonzero, so
$\deg P_{12}=1369$. The sum of the three squared residuals has exact degree
2720. The emitted live circuits have the following ledgers:
\begin{center}\small
\begin{tabular}{lrrrr}
\toprule
Presentation & Multiplications & Add/subtract & Total & Degree\\
\midrule
Three free comparisons &95&45&140&1360\\
Three explicit residuals &95&47&142&1360\\
Single product &97&47&144&1369\\
Sum of squares &98&49&147&2720\\
\bottomrule
\end{tabular}
\end{center}
The degree and coefficient in \eqref{eq:leader} are additionally checked
by a modular univariate specialization. The general leading-form argument
establishes uniformity on the actual astronomical coefficient slice.

\section{Twenty-seven unrestricted integers by an odd scale}
\label{sec:integer27}\label{sec:integer}

The ordinary and normalized fifteen-coordinate certificates admit a
smaller integer-domain conversion than the seven-positive-port conversion
in package~14. Require positivity only for
\[
 J,F,\alpha,z,\eta,\zeta,
\]
put $i=2r+1$, $s=2v+1$, and $w=u^2$, and allow
$r,v,u,f,h,T,\tau,y,\gamma$ to be arbitrary integers. Replacing each of
the six positive ports by its own copy of
\[
 \operatorname{Pos}(x,y,z)=1+x^2+y^2+z(z+1)
\]
gives $6\cdot3+9=27$ unrestricted integer unknowns. The map is positive
and onto the positive integers: represent $4n-3$ by three squares and
separate its one odd square from its two even squares.

The new point is the signed scale $s$. The source formulas still give
$q\ge16$, $C\ge0$, $q(q^2-1)<R<q^4-2$, and $|Y|\ge q^3$.
The usual unit arguments use squares and remain valid for signed roots.
The main norm excludes $X=0$: then $D_{\rm main}=ac+\gamma H$ and
$\Delta=a^2+H$ would imply $0\equiv1\pmod H$. Hence $X\ge q$.

Suppose $Y=-Z<0$. Define
\[
 B=(X+1)Z-2,\qquad c_0=-c=kZ-\eta,\qquad P_0=2XZ^2+1.
\]
Then $k(Z-1)<c_0<kZ$, $P_0>B$, and $c_0>k$. The first and main
Pell equations yield
\[
 k=2\psi_{P_0}(n),\qquad c_0=\psi_B(p),\qquad p>n,
 \qquad 2n=R+\varepsilon+v_0|E|,
\]
where $v_0\ge0$ and $\varepsilon\in\{-1,1\}$. Thus
$n\ge(R-1)/2\ge24$, $c_0>B(B^2-1)^2$, and $c_0>2p,2R$.
The normalized strong-rank argument, or its ordinary counterpart, and
the two signed auxiliary congruences give $p=R$. The sign change
$c=-c_0$ changes only a sign removed when the first congruence is squared.
The bound $|E|>R+2$ excludes wrapping. If $\varepsilon=-1$, Pell
duplication and $\chi_B(2)>P_0$ give
$c_0=\psi_B(2n+1)>Bk>kZ$, a contradiction. Therefore
$R=2n-1=2r_0+1$.

The elementary Pell growth estimates now imply
\[
 Z>\frac{c_0}{k}\ge
 \frac12\left(\frac{(2B-1)^2}{2P_0}\right)^{r_0}
 >\frac{X^{r_0}}2>2^R.
\]
The strict middle inequality is justified by
\[
 (2B-1)^2-2P_0X
 =4(2X+1)Z^2-20(X+1)Z+25-2X>0.
\]
At $Z=4$ its value is $46X+9$, and it increases thereafter; here
$X\ge16$ and $Z\ge4096$.
For the original negative Pell parameter $a+2=-B$, the two Pell projections
modulo $H=4a+3$ remain $2^R$ and $2^{-R}$. Since $R$ is odd and
$c=-\psi_B(R)$, the two signs of the main root respectively imply
\[
 X+2^R\equiv0\pmod H
 \quad\hbox{or}\quad
 X2^R\equiv1\pmod H.
\]
They contradict $0<X+2^R<|H|$ and $1<X2^R<XZ<|H|$. Thus $s>0$.
The inherited signed compiler theorem now applies. This exclusion only
needs a nonzero strong multiplier; the odd parametrization guarantees
that requirement.

Completeness follows from an exact parity identity, rather than an
additional padding conjecture. Write $q=2^t$, $M_0=q^2$, and
\[
 S_0=(C-1)+qF,\qquad
 T_0=(M_CJ+1)+q(M_F^{\rm nat}J-1),
\]
where $M_F^{\rm nat}$ is the unshifted mask. Native mask disjointness gives
$S_0\mathbin{\&}T_0=0$, and
$\operatorname{popcount}(T_0)=t+2$. Consequently
\[
 R=(M_0-1-S_0)M_0+(S_0+T_0),\qquad
 \operatorname{popcount}(R)=3t+2.
\]
The half-binomial identity gives $v_2(Y)=3t$, so $s=Y/q^3$ is odd.
The existing odd-multiplier and square-$w$ completeness choices are
unchanged. The new substitution costs $6\cdot7+2+2+1=47$ gates.

\begin{center}
\begin{tabular}{lrrr}
\hline
Construction & Equations & Degree & Gates\\
\hline
Normalized product & 1 & 274 & 114\\
Reduced normalized product & 1 & 216 & 115\\
Ordinary product & 1 & 172 & 116\\
Shifted ordinary system & 6 & 46 & 107\\
Ordinary sum of squares & 1 & 92 & 124\\
\hline
\end{tabular}
\end{center}
All rows use 27 unrestricted integers. The full signed proof, emitted
circuits, and reproducible arithmetic checks accompany the report in
\src{proofs/ODD_SCALE_27_PROOF.md}. The inherited signed and ordinary
rank arguments are those of packages 14 and 12~\cite{signed29,gap}.

The degrees are uniform over every fixed $\beta>0$. To see this without
a numerical coefficient specialization, write $p_*$ for the highest
homogeneous part of a substituted port, and set
\[
 Q_*=\beta J_*,\quad k_*=\eta_*+\zeta_*,\quad
 C_*=Q_*-F_*-\alpha_*,\quad
 N_{t,*}=w_*C_*-z_*Q_*.
\]
Here $w_*=u^2$, $s_*=2v$, and $i_*=2r$. Each guarded highest part is
a sum of three squares in its own disjoint variables. In particular
$N_{t,*}\ne0$, because its specialization at $u=0$ is $-z_*Q_*\ne0$.
The highest parts of the three product polynomials in table order are
\[
\begin{aligned}
 -8Q_*^{91}h\gamma i_*^4k_*^{13}w_*^{13}s_*^{26}T^2f^2N_{t,*},\\
 -8Q_*^{71}h\gamma i_*^2k_*^9w_*^{11}s_*^{20}T^2f^4N_{t,*},\\
  8Q_*^{55}h\gamma i_*^2k_*^9w_*^7s_*^{16}T^2f^4N_{t,*}.
\end{aligned}
\]
These are nonzero polynomials of degrees 274, 216, and 172. The unique
degree-46 system residual has highest part $a_*^2c_*^2T^2f^4$,
where $a_*=w_*s_*Q_*^4$ and $c_*=k_*s_*Q_*^3$; its square gives
the exact SOS degree 92. No value of the other fixed coefficient ports
affects these conclusions.

\section{Twenty-six unrestricted integers by robust separating weights}
\label{sec:integer26}

The twelve-coordinate norm also admits an integer-domain refinement.
Keep seven positive ports
$J,F,\alpha,z,i,\eta,\zeta$, replace them by disjoint copies of
$\operatorname{Pos}$, put $s=2v+1$ and $w=u^2$, and leave $h,T,t,v,u$
unrestricted. This uses $7\cdot3+5=26$ integers. In this construction
$i$ is positive, with no assertion that it is odd. That distinction is
necessary because the normalized multiplier is constrained to be
divisible by $H$.

Use the common prefix and write
\[
 \begin{gathered}
 P=X+ac,\quad S=1+\Delta c^2,\quad
 L=EkY,\quad A=1+L(L+k),\\
 G=Hic^2,\quad D=1+\Delta G^2,\quad
 Q=\Delta(D-1),\quad d=Q-1,\\
 v_0=cT,\quad b=c+RD,\quad
 e=Q(Dv_0^2+b^2)-1,\quad B_0=2Qv_0b.
 \end{gathered}
\]
Replace the positive-domain separating weights by
\[
 \begin{split}
 U_0&=(T^2+1)(a^2+5)H^2ic^4,\\
 M&=1+2(U_0+b),\qquad N=1+H^2M(A+1).
 \end{split}
\]
Use these weights in the same explicitly expanded two-quadratic-norm
polynomial, denoted $\mathcal N_{26}$. The formulations are
\[
 N_k=1,\quad N_t=1,\quad\mathcal N_{26}=0,
 \qquad\hbox{or}\qquad
 N_kN_t(1+\mathcal N_{26})-1=0.
\]

Here are the additional signed-domain details. The unit guard is valid
because the inner norm coefficient is even, so $\mathcal N_{26}$ is a
square modulo four. Weak transport therefore gives $C\ge0$ and
$0<R<q^4-2$ before roots are recovered. Even if $X=0$, the odd scale
ensures $|a|\ge q^3$, $|c|>1$, and
\[
 \Delta\ge(q^3-1)(q^3-3)>q^4>R,
 \qquad S>1,\quad A\ge1.
\]
Also $D>\Delta+|c|$, which gives the crucial interval
\[
 1<b<Q,
 \qquad Q-b=(\Delta-R)D-\Delta-c>0.
\]

Consider the algebraic triples satisfying
\[
 (P+H\gamma)^2=S,\qquad \tau^2=A,\qquad
 (dy^2-e)^2=B_0^2D.
\]
Every complex auxiliary branch obeys
$|y|<2(|v_0|\sqrt D+b)<M$. Indeed
$\sqrt D<|a+2|\,|G|$, and the factors in $U_0$ bound respectively
$|T|$, $|a+2|$, $|H|$, and $|c|^3$ by integer polynomials.
Furthermore
\[
 \frac{N\sqrt S}{|H|}>M(A+1)
 \ge M\sqrt A+M>M\sqrt A+|y|.
\]
Thus $(\gamma,\tau,y)\mapsto y+M\tau+N\gamma$ is injective on
distinct triples, including at $A=1$ and when the auxiliary quartic
has repeated roots. The exact norm factorization and Galois invariance
of the supplied integer $t$ imply that $\gamma,\tau,y$ are rational.

Both $\tau$ and $P+H\gamma$ are integers, so the denominator of
$\gamma$ divides $H$. For $y=r/s$ in lowest terms, the quartic
\[
 d^2y^4-2dey^2+e^2-B_0^2D=0
\]
shows that every prime dividing $s$ divides $d$. Because $Q$ is
divisible by $H^2$, $\gcd(d,H)=1$. The identity
$y+N\gamma=t-M\tau$ therefore makes $y$ integral, and
$N\equiv1\pmod H$ then makes $\gamma$ integral. This argument
does not assume that $T$ or $B_0$ is nonzero.

If $T=0$, then $dy^2=Qb^2-1$, hence
\[
 d(y^2-b^2)=b^2-1>0.
\]
Integrality gives $|y|\ge b+1$, whereas $b<Q$ gives $d\ge b$.
The left side is then at least $b(2b+1)>b^2-1$, a contradiction.
Therefore $B_0\ne0$, and
\[
 f=\frac{e-dy^2}{B_0},\qquad f^2=D
\]
is an integer. The full normalized strong and auxiliary equations have
now been recovered. The main norm excludes $X=0$ modulo $H$. The
negative-scale argument of the preceding section applies to the nonzero
normalized multiplier $Hi$ and excludes $s<0$. The inherited signed
compiler theorem then proves soundness.

For completeness, retain the canonical square $w$ and odd $s$. Choose
the strong Pell index to be a common multiple of $R$ and a period of
the invertible Pell matrix modulo $Hc^2$. This gives a positive
$i=\psi_{a+2}(m)/(Hc^2)$ and the required auxiliary quotient, without
requiring odd $i$. The canonical roots make
$t=y+M\tau+N\gamma$ an integer, and all seven positive ports have
$\operatorname{Pos}$ representations.

The exact integer degrees are
\[
 \deg d=106,\quad \deg b=\deg M=92,\quad\deg N=154,
 \quad \deg(P^2-S)=29.
\]
The two main-root branches have degrees $9$ and $-2$; therefore the
$N\gamma$ term dominates every linear factor of the norm. Its leading
form is
\[
 (\mathcal N_{26})_{3544}
 =d_{\rm top}^{8}N_{\rm top}^{16}
   (H_{\rm top}c_{\rm top}^{2})^8\ne0.
\]
The comparison system has degree 3544 and 195 gates; the product has
degree 3560 and 199 gates. A sum of squares has degree 7088 and 202
gates. All counts include the 52 substitution gates. These constructions
provide a variable-count improvement at a substantial degree and circuit
cost; the 27-integer degree-92 polynomial remains a useful alternative.

\section{A 5263-tile compiler and smaller fixed coefficients}
\label{sec:coefficients}

This section supplies the complete coefficient tuple used in
Theorem~\ref{thm:main}.  It combines two independent changes: a smaller
finite local language and a more compact placement of its copy fields.
The arithmetic equations in the preceding sections do not change.
Relative to package~11's T4 backend, the tile count falls from 37075 to
5263 and the fixed base falls from $2^{5^{18}}$ to $2^{5^{16}}$.
The general native compiler and its Pell interface are inherited from
the established blank-input construction~\cite{blank,t4,rank}; the new
work below verifies the local language and all changed layout hypotheses.

The exact machine table is included as \src{coefficients/k215.tm}, with
its literal 430-action projection in
\src{coefficients/k215_transitions.json}.  The ordered new alphabet is
\src{coefficients/windows5263.json}; its generating predicate and finite
checks are in \src{coefficients/compiler5263.py}.  These files give a
terminating finite definition of the coefficients, without running K215
until it halts.  In particular, the definition does not assume that a
ZFC refutation exists.

\subsection{The finite local language}

There are 439 cell symbols.  Codes 0 and 1 denote headless tape bits.
For $q\in\{0,\ldots,215\}$ and $b\in\{0,1\}$, code $2(q+1)+b$
denotes bit $b$ carrying a head in state $q$.  Use the five additional
symbols
\[
 L_0=434,\quad Q_0=435,\quad R_0=436,\quad V_0=437,\quad H_0=438.
\]
The first three are initialization phases, all with blank payload;
only $Q_0$ carries a head, in state 0.  The symbols $V_0,H_0$ are the
vertical and horizontal boundaries and carry no tape payload.  The
literal transition entering state 215 retains its original write and
movement.  Only after entering state 215 do we use the completion
$(215,b)\mapsto(b,0,215)$ for padding.

A tile is a tuple $(l,c,r,n)$ giving left, center, right, and next-time
center cells.  The following horizontal pair rule will be used on both
$(l,c)$ and $(c,r)$:
\begin{enumerate}
\item Two $V_0$ cells cannot be adjacent.  Immediately right of $V_0$ an
initialization phase must be $L_0$; immediately left of $V_0$ it must be
$R_0$.
\item If neither member is $V_0$, the pair agrees on being $H_0$ and
agrees on being an initialization phase.  When both are phases, the
allowed ordered pairs are $L_0L_0,L_0Q_0,Q_0R_0,R_0R_0$.
\end{enumerate}
Apart from the exceptional forbidden tile specified below, admit exactly
the tuples with these properties:
\begin{enumerate}
\item At most one of $l,c,r$ carries a head, and both horizontal pair
rules hold.
\item If $V_0$ occurs in $l,c,r$, every ordinary cell in that triple is
the blank headless cell 0, and $Q_0$ is absent.
\item The center $c$ is $V_0$ if and only if $n$ is $V_0$.  A center
$V_0$ imposes no further transition condition.
\item If $c=H_0$, then $n\in\{L_0,Q_0,R_0\}$.
\item At an interior center, a head cannot move into an adjacent
$V_0$.  If $n=H_0$, every head in the horizontal triple is in state 215.
\item Otherwise, at an interior center, $n$ is an ordinary nonphase
cell given by the exact deterministic three-neighbor update.  Its bit
is the bit written at the center, or the unchanged bit if the center
is headless.  Its head, if any, is the unique incoming head from the
left, center, or right according to the completed transition table.
Here $V_0$ is interpreted as blank and headless, and $H_0$ cannot supply
an interior update.
\end{enumerate}
The retained original one-head restriction makes the successor unambiguous.
The blank-margin restriction also excludes a head next to a vertical
boundary, so the no-escape clause is redundant for such triples; it is
stated to make the relationship to the original T4 predicate explicit.

Put the two distinguished tiles first:
\begin{equation}
 S_0=(434,435,436,1),\qquad E_0=(438,0,0,438).
 \label{eq:compiler-distinguished}
\end{equation}
The tile $S_0$ passes the displayed rules: K215's first action is
$(0,0)\mapsto(1,+1,93)$.  Admit $E_0$ as the sole exception to the
additional restrictions.  It was already a valid original T4 tile, but
it has no global extension: a tile immediately to its left would have
center $H_0$ and right neighbor 0.  Such a tile is not $E_0$ and violates
the horizontal pair rule.  Thus using $E_0$ as selector 1, which the
arithmetic forbids, discards no actual computation.

Order $S_0$ first, $E_0$ second, and all remaining passing tuples
lexicographically in the displayed coordinate order.  The finite rule,
the literal transition table, and this ordering determine the alphabet
completely.

\begin{proposition}[Exact alphabet size]
The ordered alphabet has exactly 5263 tiles.
\end{proposition}
\begin{proof}
For a headless ordinary center and no vertical boundary, the two
neighbors give
$2^2+2\cdot432\cdot2=1732$ ordinary successor tiles.  A terminal
successor allows only the two halted headed symbols, giving
$2^2+2\cdot2\cdot2=12$ further tiles.  Center~0 additionally has the
three boundary neighborhoods $(V_0,0,0),(0,0,V_0),(V_0,0,V_0)$, each
with an ordinary and a terminal successor, and the single exceptional
$E_0$.  Thus center~0 contributes 1751 tiles and center~1 contributes
1744.

Each of the 430 ordinary working-head center symbols has four choices
of its two headless neighbors and one successor, contributing 1720.
Each of the two halted-head center symbols has four ordinary and four
terminal cases, contributing 16.  The centers $L_0,R_0$ each have three
ordinary and two terminal cases; $Q_0$ has just $S_0$.  A $V_0$ center
has three admissible cells independently on either side, hence nine
cases.  An $H_0$ center has two neighbor choices on either side and
three phase successors, hence twelve cases.  Summing gives
\[
 1751+1744+1720+16+5+5+1+9+12=5263.
\]
This classification is independent of the exhaustive program.  The
program separately checks all 5215 permitted horizontal triples
against all 439 successors, or 2289385 cases, before adding $E_0$.
\end{proof}

\subsection{Why the restricted language still describes a blank halt}

Every new tile, including $E_0$, is an old T4 tile.  Therefore every new
marked periodic plane is an old marked periodic plane.  The inherited
finite-strip argument supplies a blank K215 computation ending in halt.
This is the soundness direction; restricting the alphabet cannot create
an accepting computation.

For completeness, suppose a finite blank halting computation is given.
Choose two untouched blank headless tape cells between every visited
head position and each vertical boundary.  A canonical strip has a full
$H_0$ row, then an initialization row $L_0^uQ_0R_0^v$, then the true
computation, and an arbitrarily long absorbing halted tail.  All new
horizontal conditions hold.  A row immediately before an $H_0$ row is
halted, so the terminal-neighbor condition also holds.

The inherited helical presentation changes the time coordinate by one
when crossing a vertical seam.  Within distance two of the seam, cells
are blank headless symbols, the appropriate boundary-adjacent phases,
or $H_0$.  They therefore satisfy the new restriction even when the
two sides represent different times.  The symbol $V_0$ separates them,
so their phases need not agree across the seam.  Width and height remain
arbitrarily large independent integers, and may still be chosen as
powers of five.  This presentation has one canonical Start and no
$E_0$.  Uniqueness of Start is used for this chosen completeness
witness, not assumed for arbitrary putative solutions.

\subsection{The four fixed integers}

Use a subscript $c$ for the clause radix, to distinguish it from the
polynomial $A$ in Section~\ref{sec:interface}.  Put
\begin{gather}
 A_c=8192=2^{13},\qquad \rho=2^{78125},\qquad
 B=\rho^{1953125}=2^{152587890625}=2^{5^{16}},\label{eq:newbase}\\
 M_0=7021,\qquad H_c=359389,\qquad
 T_1=738963,\qquad T_2=1118101.\notag
\end{gather}
Write $b_c=78125$ for the binary radix exponent and $L_c=1953125$
for the cell length, so $\rho=2^{b_c}$ and $B=\rho^{L_c}$.
For $0\le u\le438$, define the copy positions
\begin{align}
 e_l(u)&=6580+u,& e_c(u)&=6141+u,\notag\\
 e_r(u)&=5702+u,& e_n(u)&=5263+u.\label{eq:newcopies}
\end{align}
The support $\mathcal E$ consists of the selector positions
$0,\ldots,5262$, the four copy intervals in \eqref{eq:newcopies}, the
dummy position 7019, and the four anchors
\[
 M_0,3M_0,9M_0,27M_0=7021,21063,63189,189567.
\]
Thus $|\mathcal E|=7024$.  The possible upper dummy spill is position
7020, still strictly below the first anchor.  Set
\begin{equation}
 \mu=A_c-2+\sum_{j=1}^{2840}A_c^j.\label{eq:newmu}
\end{equation}
For the $i$th tile $w_i=(l_i,c_i,r_i,n_i)$, define
\begin{align}
 v_i={}&1+A_c^{1+l_i}+A_c^{440+c_i}+A_c^{879+r_i}
       +A_c^{1318+n_i}+\sum_{t=1}^4 A_c^{1756+t},\notag\\
 v_{e_l(u)}={}&A_c^{1+u},&v_{e_c(u)}={}&A_c^{440+u},\notag\\
 v_{e_r(u)}={}&A_c^{879+u},&v_{e_n(u)}={}&A_c^{1318+u},\notag\\
 v_{3^{t-1}M_0}={}&A_c^{1756+t}\quad(1\le t\le4),&v_{7019}={}&0.
 \label{eq:newclauses}
\end{align}
These cases have disjoint index sets.  Now put
\begin{align}
 D_C={}&\rho^{878}+\rho^{359828}+\rho^{56168}+\rho^{168504}
       +\sum_{e\in\mathcal E}v_e
          \bigl(\rho^{T_1-e}+\rho^{T_2-e}\bigr),\notag\\
 D_R={}&\rho^{H_c},\notag\\
 M_F^0={}&\mu(\rho^{T_1}+\rho^{T_2})
       +\sum_{u=0}^{438}\rho^{6141+u}
       +\rho^{63189}+\rho^{189567}\notag\\
 &+\sum_{u=0}^{438}
       \bigl(\rho^{H_c+6580+u}+\rho^{H_c+6141+u}\bigr),\notag\\
 M_C={}&B-1-\sum_{\substack{e\in\mathcal E\\e\ne1}}\rho^e
             -2\rho^{7019},\qquad M_F^{\rm nat}=M_F^0+4.
 \label{eq:newmasks}
\end{align}
The four ports used everywhere in this article are exactly
\begin{equation}
 \boxed{\quad\beta=B-1,\qquad \kappa=D_C+BD_R,\qquad
 M_C,\qquad M_F=M_F^{\rm nat}+B-1.\quad}
 \label{eq:coefficientbinding}
\end{equation}
The canonical machine-readable specification is
\src{coefficients/T5263_CONSTANTS.json}.  Its labels \src{K} and
\src{MF_source} mean $\kappa$ and $M_F$ in this article.
No individual coefficient may be substituted from another recipe.
All exponents and summation limits in
\eqref{eq:newbase}--\eqref{eq:coefficientbinding} are fixed integers.
They define ordinary integer coefficients, with no variable exponent
or extra quantified parameter.

\subsection{Population and carry conditions}

In the standard compiler notation, the new values are
\[
 \ell=13,\quad p_c=1080,\quad m_c=7023,
 \qquad m_c=2\ell+11\cdot439+8+2p_c.
\]
There is one dummy because
$m_c+1-5263-4\cdot439-4=1$.  The clause mask has population
\[
 \pc(\mu)=\ell-1+2840=2852.
\]
The two clause bands, three copy-overlap groups, and two tested anchors
are disjoint.  The extra bit~2 of $M_F^{\rm nat}$ is disjoint from all
of them.  Hence
\begin{align}
 \pc(M_F^{\rm nat})&=2\cdot2852+3\cdot439+2+1=7024,\notag\\
 \pc(M_C)&=5^{16}-7024.\label{eq:newpopulation}
\end{align}
For the low mask, omitting selector position~1 restores one otherwise
removed bit, while subtracting $2\rho^{7019}$ removes the upper dummy
bit.  Thus its missing population is exactly $7024$.  The same
disjointness and unused high positions give
\[
 0<M_C,M_F^{\rm nat}<B-1,\qquad
 M_C\equiv2\pmod4,\qquad M_F^{\rm nat}\equiv4\pmod8.
\]
This verifies the numerical portion \eqref{eq:maskcontract} of the
compiler contract.

To state the carry bound explicitly, let
$\Sigma=\sum_{e\in\mathcal E}v_e$ and
$\mathcal M=(7024+2)(2\Sigma+6)$.  Every selector coefficient has
nine $A_c$-monomials counting its constant term, while each copy and
anchor coefficient has one.  Therefore the inherited sufficient radix
exponent can be expressed using the fixed integer
\[
 U_c=4(7023+3)(18(5263+4\cdot439+4)+6)+8.
\]
It is
\begin{align*}
 b_{\min}=\max\{&13(4\cdot439+5+1080)+3,\\
 &13(4\cdot439+4)+\operatorname{bitlength}(U_c)+1,\ 4\}\\
 ={}&36936.
\end{align*}
Our exponent $78125$ exceeds it.  An exact integer calculation also
gives
\[
 \operatorname{bitlength}(\mu)=36921,\qquad
 \operatorname{bitlength}(4\mathcal M+8)=22909.
\]
Thus $\rho>2\mu+4$ and $\rho>4\mathcal M+8$, with ample room.  These
are the original sufficient carry bounds, including rotations and the
optional upper dummy bit; a comparison only of the displayed scalar
exponents would not have established that fact.

\subsection{Convolution fields and synchronization}

At target $e_c(u)$, the bare shift $878=2\cdot439$ has exactly the
source $e_n(u)$.  At the targets $H_c+e_l(u)$ and $H_c+e_c(u)$, the bare
shift $H_c+439$ has sources $e_c(u)$ and $e_r(u)$; multiplication by
$D_R$ has sources $e_l(u)$ and $e_c(u)$.  The two other bare shifts,
$8M_0$ and $24M_0$, have precisely the anchor source/target pairs
$(M_0,9M_0)$ and $(3M_0,27M_0)$.  Exact finite set calculations verify
all 1319 targets and exclude every unintended hit there.

Set $E_{\max}=27M_0=189567$.  Every copy or anchor target is below
$T_1-E_{\max}$.  At each clause band $T_j$, no bare term contributes
and the other clause band cannot contribute.  Consequently the same
clause tests select either an empty cell or one tile together with its
four matching copies and four anchors.  The field tests yield exactly
\[
 w_i(r)=w_{i-1}(c),\qquad w_i(c)=w_{i-1}(l),\qquad
 w_i(n)=w_{i-h_0}(c).
\]
Writing $a_i=w_i(c)$, they imply
$w_i=(a_{i+1},a_i,a_{i-1},a_{i-h_0})$.  Thus cyclic data pull back to
the actual plane $a_{-x-h_0t}$ for arbitrary stride $h_0$.  No
coprimality or divisibility assumption on that stride is introduced.

The synchronization check must cover arbitrary bit rotations.  Here
\begin{gather}
 T_2-T_1=379138=2E_{\max}+4\equiv2\pmod4,\notag\\
 T_2+2E_{\max}+2=1497237<1953125.\label{eq:newseparation}
\end{gather}
The support, including the upper dummy spill, lies in
$[0,E_{\max}]$.  The two possible clause-band contamination sets are
therefore disjoint modulo the cell length.  The inherited clean-band
argument always has an uncontaminated band available.

For a binary rotation $b_ct+\epsilon$, $0\le\epsilon<b_c$, the carry bound
shows that a nonzero remainder cannot satisfy both odd anchor tests:
for $1\le\epsilon\le b_c-2$ the rotated digits are even, and for
$\epsilon=b_c-1$ only the single upper dummy spill can be odd.  Therefore
$\epsilon=0$.  For a radix-aligned rotation, the two tests require a
support pair separated by $18M_0$.  The only such pair is
$(9M_0,27M_0)$, because all nonanchor support is below $M_0$ and the
other anchor differences are distinct.  The circumference is too large
to introduce a second pair by wrapping.  Hence the radix rotation is
zero modulo the cell length.  These are precisely the two original
synchronization steps, with their hypotheses rechecked for the new
support.

\subsection{Five-adic control of the actual packed rank}

The change of copy spacing also changes the residue used in the
dummy-bit adjustment, so it must be recalculated.  Since $78125$ is a
power of five, $\rho\equiv2\pmod5$.  By
\eqref{eq:newseparation}, each paired clause contribution cancels
modulo five, independently of its coefficient $v_e$.  The four bare
exponents in $D_C$ are respectively $2,0,0,0$ modulo four.  Hence
\begin{equation}
 D_C\equiv D_R\equiv2\pmod5,
 \qquad 2D_C-D_R\equiv2\pmod5.\label{eq:newunit}
\end{equation}
The former residue 3 is not required; the relevant property is
invertibility modulo five.

For clarity, use $w_0$ for the spatial width, $H_0^{\rm time}$ for the
temporal height, and $n_0=w_0H_0^{\rm time}$ for their product, to
distinguish these quantities from the polynomial coordinates.  Choose
all three as powers of five, with $H_0^{\rm time}\ge25$ and
$n_0>25d$, where $d=5^{16}$.  The canonical packed base is
$q=B^{n_0}$.  Switching on the free low dummy bit in cell $j$ changes
the \emph{actual} packed rank by $-\Gamma B^j$, where
\begin{equation}
 \Gamma=\rho^{7019}(q^2-1)
       \bigl[1+q(D_C+BD_R+B^{w_0})\bigr].\label{eq:newgamma}
\end{equation}
This formula has the original no-wrap conditions $j+1,j+w_0<n_0$.
Because all relevant exponents are powers of five,
$B\equiv q\equiv B^{w_0}\equiv2\pmod5$.  The bracket in
\eqref{eq:newgamma} is therefore
\[
 1+2(D_C+2D_R+2)\equiv2D_C-D_R\equiv2\pmod5.
\]
Also $\rho^{7019}\equiv3$ and $q^2-1\equiv3$ modulo five, so
\begin{equation}
 \Gamma\equiv3\pmod5,\qquad \gcd(\Gamma,dn_0)=1.
 \label{eq:gammatuning}
\end{equation}
The unchanged Boolean subset-adjustment theorem on the cells
$\{4j:0\le j<n_0/5\}\cup\{1\}$ therefore supplies the required
congruence $R\equiv dw_0\pmod{dn_0}$.  Width and height padding gives
all size and no-wrap conditions.  This is a congruence adjustment to
the packed expression \eqref{eq:rank}, not the selection of an
independent formal exponent.

The canonical origin test still sees Start at selector 0 and the
forbidden tile at selector 1.  Every copy and dummy position is
positive, so their new placement does not change the unit selector
contribution at the origin.  Dummy changes preserve all tile choices
and overlaps.  The original positive-slack argument and carry bounds
complete the arithmetic witnesses.  The odd exponent $d=5^{16}$ and
the canonical rank congruence $R\equiv3\pmod4$ are preserved; these
are the padding and parity properties used in the integer construction
of Section~\ref{sec:integer}.

\begin{proposition}[Complete replacement of the coefficient tuple]
The four integers in \eqref{eq:coefficientbinding} satisfy the full
blank-input compiler contract used in Section~\ref{sec:interface}.
In particular, substituting this whole tuple into the new arithmetic
circuits preserves their equivalence with a blank K215 halt.
\end{proposition}
\begin{proof}
The local language has the required marked-tableau soundness and
padded-helical completeness.  Equations
\eqref{eq:newpopulation}--\eqref{eq:newseparation} and the exact target
checks verify the numerical, carry, field, and synchronization
hypotheses of the inherited native compiler.  Equations
\eqref{eq:newgamma}--\eqref{eq:gammatuning} verify the actual-rank
tuning hypothesis.  The original Start-origin and slack arguments are
unchanged.  The blank-input compiler theorem and the arithmetic
equivalences therefore apply to this entire tuple~\cite{blank,t4,rank}.
\end{proof}

\subsection{Size and evidence boundary}

The coefficient $\beta$ has exactly $152587890625$ bits.  The largest
port is $\kappa$, with exactly
\[
 d+78125H_c+1=180665156251
\]
bits, because $D_C<B$ while $BD_R$ has its highest bit at
$d+78125H_c$.  Thus even this improved tuple has enormous expanded
numerals.  The supplied finite recipes are exact; the expansions were
not materialized.

The program recomputes the complete alphabet, verifies its inclusion
in the original T4 alphabet, checks the impossibility of the forbidden
tile through every adjacent candidate, and verifies all carry and
support conditions above.  It also checks 18 finite control tableaux
containing 927 windows, including helical seams.  Those controls use
specified altered transitions to create finite examples and do not
assert a halting run of literal K215.  The archive records the finite
checks separately from the universal mathematical proof.

The arithmetic, finite combinatorics, recurrence arguments, and
coefficient definitions used here are formalizable in ZFC.  The
general compiler and Pell results remain imported mathematical
dependencies.  Neither these finite checks nor inspection of the
archived endpoint reports constitutes a new end-to-end Lean or Rocq
verification of the Diophantine equation.

\section{Logical meaning, verification, and reproduction}\label{sec:logic}

\subsection{The equivalence proved in ZFC}

The dependency chain is
\[
 \text{integer solution}
 \ \longleftrightarrow\ \text{finite accepted blank tableau}
 \ \longleftrightarrow\ \operatorname{Halt}(K215,\mathrm{blank})
 \ \longleftrightarrow\ \ZFCRef.
\]
The first equivalence consists of the root reconstruction, the signed-domain
arguments, the inherited Pell compiler, and the finite coefficient adapter.
Completeness runs through the chain in the opposite direction. All integer
exponents in the coefficient recipe denote fixed integers; none is a hidden
exponentiation in an unknown. Each displayed final object is consequently
an ordinary polynomial, or a finite system of ordinary polynomials, over
$\Z$.

The final endpoint is the theorem
\begin{quote}\small\ttfamily
blank\_halts\_iff\_zfc\_refutation :\\
BlankHalts machine215 $\leftrightarrow$ BProv ZFCax\_s [] Form.fBot
\end{quote}
in \src{KEndpoint215.lean}. It has no consistency hypothesis or remaining
machine-specific theorem parameter. The Rocq development gives the
corresponding equivalence with its \src{StandardZFC} theory~\cite{endpoint}.

For precision, \src{BProv} means existence of a finite list of background
axioms and a finitary first-order derivation from them. The letter B refers
to the background theory; it does not impose a numerical bound on proof
length. The namespace \src{BoundedZFCConsistency} in some imported modules
does not change this endpoint into a bounded-consistency statement.
The axiom predicate includes all instances of Separation and Replacement
and a standard selector formulation of Choice. We may define
$\Con(\mathrm{ZFC})$ using this exact effective presentation. The usual
finite proof translations between standard equivalent presentations give
the same result for another conventional coding~\cite{endpoint,t4}.

The machine calibration uses finitary syntax, first-order completeness,
constructible-universe arguments, and a well-founded quotient-model
construction. These are ordinary mathematical constructions available in
ZFC; they do not require assuming $\Con(\mathrm{ZFC})$. The new algebraic
steps use finite extension fields, automorphisms of a finite normal closure,
elementary Pell arithmetic, and finite computation. They too are
formalizable in ZFC. Composing those proofs gives Theorem~\ref{thm:main},
and classical negation gives its unsolvability version.

This formulation deliberately separates two claims. We supply an ordinary
mathematical proof of a theorem provable in ZFC. We do not supply a new
machine-checked derivation in an encoded ZFC proof calculus, or a new
end-to-end Lean/Rocq formalization of the Diophantine reductions. The
external proof-assistant endpoint by itself is not such a translation.
No stronger internalization in PA is asserted here.

If ZFC is consistent, the displayed polynomial has no solution in the
standard positive integers, or in the standard integers for its signed
variants. Under the usual hypotheses of the second incompleteness theorem,
ZFC then cannot prove this lack of solutions: the proved equivalence would
turn such a proof into a proof of its own consistency. The construction
does not determine whether a solution exists. Nor do we infer from mere
consistency that ZFC cannot prove the opposite assertion; that stronger
independence conclusion would require an additional argument or hypothesis.

\subsection{What was checked, and what was imported}

The source snapshot is pinned to the full Git commit in the abstract.
Files in \src{upstream/} preserve the inherited texts and their original
historical snapshot notices. \src{SOURCE_MANIFEST.json} records their
retrieval paths, Git blob identities, and SHA-256 hashes. The literal
430-transition machine table is included and compared with the exported
transition list before the finite compiler is regenerated.

The archive also preserves two \emph{inherited} endpoint audit logs.
The Lean log reports only \src{propext}, \src{Classical.choice}, and
\src{Quot.sound}; the Rocq log reports closure under the global context.
We inspected these records and the endpoint source. We did not execute a
fresh Lean or Rocq build in preparing this archive. Their limited size
does not make them a substitute for the complete upstream developments.

The new executable checks are standard-library Python. They verify:
\begin{enumerate}
\item Circuit register discipline, liveness, exact operation counts,
  consistency with independently evaluated norm formulas, and specialized
  polynomial degree certificates. Structural leading-form arguments in
  the proofs establish the uniform exact degrees for the authentic fixed
  coefficient tuple.
\item The signed substitutions, the negative-parameter Pell projection
  identities, the population identity used to make $s$ odd, and explicit
  small negative-scale Pell tuples with both main-root signs excluded.
\item The 26-integer norm by a separate implementation of multiplication
  and conjugation in $\Z[r,u]/(r^2-A,u^2-S)$. This checks the expanded
  circuit without copying its scalar norm formulas, including $T=0$,
  $X=0$, and negative-$s$ cases.
\item Every candidate successor for the refined admissible horizontal
  triples, the exact 5263-tile enumeration, clause/copy geometry, carry
  bounds, synchronization anchors, mask populations, and the modular unit
  needed to tune the actual rank. Eighteen finite control tableaux check
  927 seam and boundary windows.
\end{enumerate}

Finite tests support the arithmetic specification and catch implementation
errors. They do not establish the universal Pell lemmas, the unbounded
tableau theorem, or consistency. In particular, the synthetic exact-root
fixtures are tests of the algebraic norm identity; they are not claimed
solutions for the K215 coefficients. The Galois reconstruction and signed
zero-case arguments are proved in the text and independent audit notes,
not inferred from testing a sample of integers.

\subsection{Files and reproducible commands}

From the extracted archive root, run
\begin{quote}\ttfamily
python3 reproduce.py\\
latexmk -pdf -interaction=nonstopmode -halt-on-error report.tex
\end{quote}
The first command needs Python 3.10 or later and its standard library only.
It builds the new arithmetic packets and the finite compiler in a temporary
directory, compares them with the bundled artifacts, checks the literal
machine binding, runs the separate quotient-algebra audit, and writes a
verification receipt. It does not expand the enormous fixed coefficients.
The second command uses a standard TeX installation to rebuild this PDF.
The corresponding \src{Makefile} provides \src{verify}, \src{pdf}, and
\src{clean} targets.

\begin{center}\small
\begin{tabular}{>{\raggedright\arraybackslash}p{4.6cm}p{8.6cm}}
\toprule
Path & Purpose\\
\midrule
\src{report.tex}, \src{sections/} & Complete article sources.\\
\src{data/polynomial12.json} & Fully substituted 12-positive-variable product.\\
\src{data/system12.json} & Its three comparisons.\\
\src{data/integer26/} & 26-integer product, system, and sum of squares.\\
\src{data/integer27/} & Five lower-degree 27-integer alternatives.\\
\src{code/} & Deterministic circuit builders and independent norm audit.\\
\src{coefficients/} & Literal machine, 5263 tiles, coefficient recipe,
  finite verifier, and inherited endpoint evidence.\\
\src{COEFFICIENT_BINDING.json} & Exact mapping from each circuit's port names
  to the one complete coefficient tuple.\\
\src{proofs/}, \src{upstream/} & New proof/audit notes and inherited proof dependencies.\\
\src{comparison.csv} & Machine-readable comparison of the stated alternatives.\\
\src{SOURCE_MANIFEST.json} & Pinned source provenance.\\
\src{MANIFEST.sha256} & SHA-256 checksums for the distributed files.\\
\bottomrule
\end{tabular}
\end{center}

\subsection{Which form to use}

For the fewest positive unknowns in this reviewed family, use
\src{polynomial12.json}; its three-equation system is slightly shorter to
evaluate. For unrestricted integer unknowns, the new 26-integer equation
uses the fewest coordinates among the explicitly bound variants presented
here. If degree or arithmetic size matters more, the 27-integer ordinary
system has degree 46 and 107 operations, while its single ordinary product
has degree 172 and 116 operations. The unchanged 15-positive-variable
product remains the shortest known evaluator in this repository comparison,
at 67 operations. A degree-four target still favors the existing quartic
lift. The refined coefficient tuple can be used coherently with any version
whose complete blank-compiler interface is proved.

The new reductions therefore add points to a tradeoff table. They do not
make one equation uniformly simplest. Further useful work would be to
formalize the separating-map and denominator lemmas, reduce the large
weights while retaining their signed bounds, or compile the smaller
numerical proof calculus directly. Those directions are not treated as
results of this report.


\begin{thebibliography}{20}\small
\bibitem{catalogue} V. Reshetnikov, ProveIt contributors.
\emph{Con(ZFC) in Diophantine form}, wave-2 reconciliation and audit.
Snapshot \src{f0ae4e0991ca}, 22 October 2026.
\href{https://github.com/VladimirReshetnikov/ProveIt/blob/f0ae4e0991ca44a7bb4e5671d42e33d775e48035/SetTheory/Cardinals/docs/reports/hilbert-tenth-problem/con-zfc-diophantine-forms/README.md}{Pinned collection README}.
\bibitem{blank} ProveIt, wave-1 package 06.
\emph{The 67-operation blank-input polynomial} and
\emph{Blank helical proof}. Authenticated copies in this archive:
\src{upstream/normalized_blank_proof.md} and \src{upstream/blank_helical_proof.md}.
\bibitem{rank} ProveIt, native-window/Pell development.
\emph{Pell lemma audit}, ordinary strong-rank interface, and signed auxiliary
congruence proofs. Copies in \src{upstream/}; exact source paths and hashes
in \src{SOURCE_MANIFEST.json}.
\bibitem{root13} ProveIt, package 09.
\emph{Smaller Diophantine descriptions of ZFC consistency}; root amalgamation
and the 13-variable weighted construction, 22 October 2026.
\href{https://github.com/VladimirReshetnikov/ProveIt/tree/f0ae4e0991ca44a7bb4e5671d42e33d775e48035/SetTheory/Cardinals/docs/reports/hilbert-tenth-problem/con-zfc-diophantine-forms/09-combined13}{Pinned package 09}.
\bibitem{t4} ProveIt, package 11.
\emph{T4 bridge audit} and exact coefficient recipe, 22 October 2026.
\href{https://github.com/VladimirReshetnikov/ProveIt/tree/f0ae4e0991ca44a7bb4e5671d42e33d775e48035/SetTheory/Cardinals/docs/reports/hilbert-tenth-problem/con-zfc-diophantine-forms/11-t4}{Pinned package 11}.
\bibitem{gap} ProveIt, package 12.
\emph{Smaller Diophantine forms of ZFC consistency}, ordinary strong and gap
constructions, 22 October 2026.
\href{https://github.com/VladimirReshetnikov/ProveIt/tree/f0ae4e0991ca44a7bb4e5671d42e33d775e48035/SetTheory/Cardinals/docs/reports/hilbert-tenth-problem/con-zfc-diophantine-forms/12-gap-norm14}{Pinned package 12}.
\bibitem{signed29} ProveIt, package 14.
\emph{Signed-domain proof} and 29-integer construction, 22 October 2026.
\href{https://github.com/VladimirReshetnikov/ProveIt/blob/f0ae4e0991ca44a7bb4e5671d42e33d775e48035/SetTheory/Cardinals/docs/reports/hilbert-tenth-problem/con-zfc-diophantine-forms/14-integer29/proofs/signed_domain_proof.md}{Pinned signed-domain proof}.
\bibitem{endpoint} ProveIt.
\emph{KEndpoint215.lean} and the K215 verification records.
\href{https://github.com/VladimirReshetnikov/ProveIt/tree/f0ae4e0991ca44a7bb4e5671d42e33d775e48035/Computability/BusyBeaver/Research/ZFC/k225-formal/verification/k215}{Pinned verification directory}.
\bibitem{mr75} Yu. Matiyasevich and J. Robinson.
Reduction of an arbitrary Diophantine equation to one in 13 unknowns.
\emph{Acta Arithmetica} 27 (1975), 521--553.
\href{https://doi.org/10.4064/aa-27-1-521-553}{doi:10.4064/aa-27-1-521-553}.
\bibitem{jones82} J. P. Jones.
Universal Diophantine equation.
\emph{Journal of Symbolic Logic} 47(3) (1982), 549--571.
\href{https://www.cambridge.org/core/journals/journal-of-symbolic-logic/article/abs/universal-diophantine-equation/113EBF3D928DFC6011C52BECAE81A087}{Publisher record}.
\bibitem{bayer} J. Bayer, M. David, M. Hassler, Yu. Matiyasevich, D. Schleicher.
Diophantine equations over $\Z$: universal bounds and parallel formalization.
\href{https://arxiv.org/abs/2506.20909v2}{arXiv:2506.20909v2}, 2025.
See also the authors' \href{https://isa-afp.org/entries/Diophantine_Universal_Pairs.html}{Archive of Formal Proofs development}.
\end{thebibliography}
\end{document}
